\documentclass[11pt,a4paper]{article}
% ============================================================
% Paper 1 - VERSI JICTRA (LaTeX, untuk cek tampilan PDF).
% Konten IDENTIK dengan paper1-jictra.txt (yg dipakai copy-paste ke template DOCX).
% Format JICTRA: Inggris, SI units, abstrak 100-200 kata, referensi gaya nomor,
%   urutan Title/Authors/Abstract/Body/Acknowledgements/Nomenclature/References,
%   target <= 15 halaman. Semua ANGKA dari eksperimen nyata (tervalidasi).
% ============================================================
\usepackage[utf8]{inputenc}
\usepackage[T1]{fontenc}
\usepackage{amsmath,amssymb}
\usepackage{booktabs}
\usepackage{graphicx}
\usepackage{geometry}
\usepackage{times}
\usepackage[hidelinks]{hyperref}
\geometry{margin=2.5cm}
\setlength{\parskip}{0.4em}
\linespread{1.0}

\title{\textbf{Semantic Feature Mapping Validated on Live Cloud Traffic:
A Deployment-Oriented Recipe for Cross-Network Intrusion Detection}}

\author{Hero Yudo Martono\thanks{Corresponding author: \texttt{hero@pens.ac.id}},
\quad Iwan Syarif, \quad Ferry Astika Saputra \\[0.3em]
\normalsize Department of Informatics and Computer Engineering, \\
\normalsize Politeknik Elektronika Negeri Surabaya (PENS), Surabaya, Indonesia \\
\normalsize \texttt{hero@pens.ac.id; iwanarif@pens.ac.id; ferryas@pens.ac.id}}

\date{}

\begin{document}
\maketitle

\begin{abstract}
\noindent Machine-learning-based network intrusion detection systems (NIDS) often
perform well within a training network but fail to generalize to other networks. One
overlooked challenge is \emph{feature-extractor mismatch}, in which different
flow-export tools compute nominally equivalent features differently. This study
proposes a statistically audited \emph{Semantic Feature Mapping} (SFM) that combines
semantic verification with scale-distribution checks and rejects incompatible feature
pairs, including TCP-window features. Using CSE-CIC-IDS2018 and UNSW-NB15, we evaluate
cross-network generalization with a lightweight XGBoost model. Although in-domain
performance reaches Matthews correlation coefficients (MCC) of 0.91 and 0.74, single
-source cross-network performance collapses to approximately zero. Controlled
experiments show the shared features retain discriminative information under joint
training, while a domain-classifier accuracy of about 0.99 evidences substantial
distribution shift. We further evaluate the approach on real AWS traffic---a field
validation step largely absent from prior cross-dataset studies. Zero-shot transfer
performs poorly, whereas calibration with 1\% of target labels raises MCC to 0.49, and
10--20\% to about 0.59, comparable to the AWS-only reference (0.598). AWS gains remain
dependent on the evaluation protocol.

\medskip
\noindent\textbf{Keywords:} network intrusion detection; cross-network
generalization; semantic feature mapping; feature-extractor mismatch; real-traffic
validation; few-shot calibration; edge deployment.
\end{abstract}

\section{Introduction}
Cyberattacks impose major and growing economic losses \cite{cremer}, making
ML-based NIDS a critical line of defense \cite{khraisat,buczak}. Such NIDS routinely
report over 99\% accuracy on public benchmarks, yet this rarely materializes as a
working detector on a different, previously unseen network. The gap stems from a
seldom-tested assumption: that training and operational traffic share the same
distribution. Prior cross-dataset studies \cite{cantone,dhooge,verkerken} confirm
the collapse---average MCC falling from 0.946 within-dataset to 0.294 cross-dataset
\cite{cantone}---but they generally stop at observing it, and they evaluate only
dataset-to-dataset, never against live production traffic.

This paper closes two deployment-facing gaps that, to our knowledge, no prior work
addresses together:
\begin{enumerate}
  \item \textbf{Feature-extractor mismatch.} Different flow exporters compute the
  same nominal feature differently; a name-based unified schema silently admits
  non-comparable pairs. We detect and reject such pairs by statistical validation.
  \item \textbf{No validation on a real new network.} Cross-dataset papers measure
  dataset-to-dataset drops; none deploys the mapped model on live traffic to measure
  false-alarm behavior and repair cost.
\end{enumerate}

\paragraph{Core novelty.} The combination absent from prior work: a statistically
validated cross-tool Semantic Feature Mapping carried \emph{end-to-end to a live
cloud network (AWS)}---measuring false-alarm rate under benign traffic, quantifying
a three-domain distribution shift, and restoring detection with a few local labels.
Unified-feature-set work \cite{sarhan} standardizes schemas but neither validates
each pair statistically nor tests on live traffic; cross-dataset studies
\cite{cantone,dhooge,verkerken} document the drop but offer no field validation or
low-cost repair. We bridge both.

\paragraph{Contributions.}
\begin{enumerate}
  \item A cross-tool SFM with per-pair statistical validation, including a concrete
  feature-extractor mismatch finding (TCP-window features).
  \item Quantitative evidence of cross-network collapse, and a controlled diagnosis
  separating feature insufficiency from distribution shift.
  \item A low-cost repair (1\% target labels, or cross-dataset mixup) with
  multi-seed confidence intervals.
  \item End-to-end validation on live AWS traffic: false-alarm rate, a quantified
  three-domain shift, and few-shot recovery---the piece missing from prior
  cross-dataset work.
\end{enumerate}

\paragraph{Research questions.} RQ1: can semantically equivalent flow features from
heterogeneous tools be reliably identified? RQ2: how far does a single-source NIDS
generalize after semantic alignment? RQ3: what distinguishes representation
insufficiency from distribution shift? RQ4: how much target labeling recovers
performance? RQ5: does the recipe transfer to real AWS traffic?

\paragraph{Reproducibility.} All reported values come from executed experiments; no
illustrative values are presented as results, and negative/non-monotone outcomes are
reported as observed.

\section{Related Work}
ML-based NIDS span tree ensembles and deep models (e.g., deep and transformer-based
detectors \cite{shone,wu}, and XGBoost variants on UNSW-NB15 \cite{husain}), most
reporting strong in-domain results that do not guarantee cross-network robustness.

\paragraph{Cross-dataset generalization.} D'Hooge et al.\ \cite{dhooge} showed
supervised NIDS fail on unseen networks; Verkerken et al.\ \cite{verkerken} reported
about 30\% AUROC drop with unsupervised models; Cantone et al.\ \cite{cantone}
documented MCC falling from 0.946 (within) to 0.294 (cross) over four datasets.
These works confirm the collapse but stop at observing it, without separating
feature insufficiency from distribution shift and without a low-cost recovery or
live-traffic validation.

\paragraph{Unified feature sets.} Sarhan et al.\ \cite{sarhan,sarhannf} proposed a
standardized NetFlow feature set to enable cross-dataset comparison. We go further:
each mapped pair is statistically validated on its original scale, and spuriously
named matches are rejected (the TCP-window case)---exactly what a name-based schema
admits silently.

\paragraph{Dataset integrity.} Reconstructions found over 20\% of CIC-IDS2017 flows
needed relabeling \cite{engelen}, CICFlowMeter feature errors \cite{rosay,lanvin},
and 6--7\% label-error rates on CIC/CSE-CIC datasets \cite{liu}. Catillo et al.\
\cite{catillo} questioned whether benchmark progress reflects real advances. Layeghy
et al.\ \cite{layeghy} showed synthetic benign traffic differs statistically (large
Wasserstein distance) from real networks---motivating both our cross-tool mapping
and our live-traffic validation.

\paragraph{Domain adaptation.} CORAL \cite{coral} aligns covariance; few-shot methods
target label-limited NIDS. Deep domain-invariant NIDS learns adversarially invariant
features but needs a heavy deep stack and ample target data---impractical for edge.
We keep a fixed lightweight XGBoost \cite{xgboost} backbone and recover performance
with cheap per-deployment calibration, a different edge-oriented operating point. A
head-to-head against deep DA needs a different stack and is left as future work.

\section{Data and Semantic Feature Mapping}
\paragraph{Datasets (deliberately cross-source).} CSE-CIC-IDS2018 \cite{cicids}
(CICFlowMeter, 68 clean features) and UNSW-NB15 \cite{unsw,unsweval} (Argus and
Bro/Zeek, 42 features; 175{,}341 train and 82{,}332 test records). The differing
extraction tools are a prerequisite for realistically testing feature-extractor
mismatch, not a weakness. All attack types are collapsed to a single
\texttt{attack} class (controlled binary setting).

\paragraph{SFM construction and validation.} SFM maps functionally equivalent
features even when names differ (e.g., \texttt{Flow Duration}$\leftrightarrow$\texttt{dur}).
Each candidate pair is not accepted on name similarity but validated in three steps:
(i) match semantics from each tool's official documentation; (ii) return CIC features
to original scale via Eq.~\eqref{eq:unscale} (our CIC release is z-scored); (iii)
compare summary statistics (median, p99, max) for scale-order equivalence. A pair is
\emph{aligned} when median/p99 ratios share an order of magnitude,
\emph{scale-mismatch} at 1--3 orders (unit conversion), \emph{different-feature}
beyond 3 orders. The nine strong pairs form Model~A; two inter-arrival-time pairs
form optional Model~B; two TCP-window pairs are discarded as a feature-extractor
mismatch (continuous in CICFlowMeter, effectively binary in Argus/Bro) despite
similar names. Per-dataset z-score normalization resolves residual unit differences.

\begin{equation}
  x_{\mathrm{orig}} = x_{\mathrm{scaled}} \cdot \sigma + \mu
  \label{eq:unscale}
\end{equation}

\begin{table}[htbp]
\centering
\caption{Semantic Feature Mapping (SFM): CSE-CIC-IDS2018 $\leftrightarrow$ UNSW-NB15,
validated on the original scale. Medians on original scale (CIC un-scaled first).}
\label{tab:sfm}
\footnotesize
\setlength{\tabcolsep}{4pt}
\begin{tabular}{llrrl}
\toprule
CIC feature & UNSW feature & CIC med. & UNSW med. & Verdict \\
\midrule
Flow Duration    & dur    & $20761$ & $0.0016$ & Model A ($\mu$s vs s) \\
Tot Fwd Pkts     & spkts  & $2$ & $2$ & Model A; aligned \\
Tot Bwd Pkts     & dpkts  & $1$ & $2$ & Model A; same order \\
TotLen Fwd Pkts  & sbytes & $43$ & $430$ & Model A; comparable \\
TotLen Bwd Pkts  & dbytes & $101$ & $164$ & Model A; aligned \\
Fwd Pkt Len Mean & smean  & $36$ & $73$ & Model A; aligned \\
Bwd Pkt Len Mean & dmean  & $66$ & $44$ & Model A; aligned \\
Flow Byts/s      & sload  & $787.5$ & $8.8{\times}10^{5}$ & Model A; byte/s vs bit/s \\
Bwd Pkts/s       & dload  & $3.70$ & $1447$ & Model A; pkt/s vs bit/s \\
\midrule
Fwd IAT Mean     & sinpkt & $3544$ & $0.28$ & Model B ($\mu$s vs ms) \\
Bwd IAT Mean     & dinpkt & $0$ & $0.006$ & Model B; optional \\
\midrule
Init Fwd Win Byts & swin  & $953$ & $0$ & \textbf{Discarded}: binary vs cont. \\
Init Bwd Win Byts & dwin  & $-1$ & $0$ & \textbf{Discarded}: binary vs cont. \\
\bottomrule
\end{tabular}
\end{table}

\paragraph{Key finding (TCP-window mismatch).} UNSW \texttt{swin}/\texttt{dwin} are
practically binary ($\{0,255\}$) while CIC \texttt{Init Fwd/Bwd Win Byts} are
continuous ($0$--$65535$). Despite both being named ``TCP window,'' the data are not
comparable, so we discard the pair and report it as a feature-extractor mismatch.
Final Model~A (9 features): \texttt{dur, spkts, dpkts, sbytes, dbytes, smean, dmean,
sload, dload}; z-score is fit on train only (no leakage).

\section{Cross-Network Baseline and Diagnosis}
\paragraph{Collapse (RQ2).} A single XGBoost excels in-domain (MCC 0.91 on CIC, 0.74
on UNSW) but collapses cross-network (MCC $\approx 0$)---consistent with the
literature but here measured on the SFM-aligned feature space, ruling out a trivial
schema mismatch as the cause. We report MCC as the primary metric
(Eq.~\eqref{eq:mcc}) because it is robust under heavy class imbalance.

\begin{table}[htbp]
\centering
\caption{Cross-dataset baseline (MCC, Model~A). High in-domain performance collapses
cross-network; cross-network values are indistinguishable from zero over 5 seeds.}
\label{tab:baseline}
\begin{tabular}{lcc}
\toprule
Scenario & Model A & Model B \\
\midrule
Same-CIC (train \& test CIC)   & $0.913$ & $0.914$ \\
Same-UNSW (train \& test UNSW) & $0.745$ & $0.743$ \\
CIC$\rightarrow$UNSW           & $-0.072$ & $0.004$ \\
UNSW$\rightarrow$CIC           & $-0.061$ & $-0.046$ \\
\midrule
\emph{Generalization gap}      & $0.81$--$0.99$ & $0.79$--$0.91$ \\
\bottomrule
\end{tabular}
\end{table}

\begin{equation}
  \mathrm{MCC} = \frac{TP\cdot TN - FP\cdot FN}
  {\sqrt{(TP{+}FP)(TP{+}FN)(TN{+}FP)(TN{+}FN)}}
  \label{eq:mcc}
\end{equation}

\paragraph{Diagnosis: shift, not (only) insufficiency (RQ3).} Joint training on the
nine SFM features performs well on both domains, showing the shared features carry
sufficient discriminative information---so feature insufficiency is not the sole
explanation. A domain classifier separates CIC from UNSW with $\approx 0.99$ accuracy
and a large inter-domain 1-Wasserstein distance, directly measuring the distribution
shift as the major contributor. This is conceptually motivated by the Ben-David bound
on target-domain error (Eq.~\eqref{eq:bendavid}), where a larger domain divergence
upper-bounds achievable cross-domain performance.

\begin{equation}
  \epsilon_T(h) \le \epsilon_S(h) + d_{\mathcal{H}\Delta\mathcal{H}}(S,T) + \lambda
  \label{eq:bendavid}
\end{equation}

\section{Low-Cost Recovery (RQ4)}
Minimal calibration recovers cross-network detection: using only 1\% of
target-network labels, or supervised cross-dataset \emph{mixup} that uses only target
train labels, restores cross-network MCC into the 0.65--0.91 range (multi-seed, with
confidence intervals). An adaptation-direction asymmetry is observed and analyzed:
CIC$\rightarrow$UNSW recovers lower (about 0.66--0.69) than UNSW$\rightarrow$CIC
(about 0.90), explained by greater domain dispersion in UNSW (mean SFM-feature
variance 3.64 vs 0.68 for CIC, about $5.3\times$), concentrated in rate/packet-size
features. We report this as an evidence-consistent interpretation, not a causal proof.

\begin{table}[htbp]
\centering
\caption{Few-shot target adaptation (cross-network MCC; mean over 5 seeds). Mixup
reported on a single seed (42).}
\label{tab:fewshot}
\begin{tabular}{lcc}
\toprule
Target label fraction & CIC$\rightarrow$UNSW & UNSW$\rightarrow$CIC \\
\midrule
$0\%$ (baseline) & $-0.074$ & $-0.094$ \\
$1\%$            & $0.659$  & $0.900$ \\
$5\%$            & $0.689$  & $0.909$ \\
$10\%$           & $0.691$  & $0.911$ \\
\midrule
Mixup (target train labels) & $0.680$ & $0.798$ \\
\bottomrule
\end{tabular}
\end{table}

\section{Real-Traffic Validation on AWS (RQ5)}
We deploy the SFM-aligned detector on live AWS cloud traffic, the step prior
cross-dataset work omits. Three consistent findings:
\begin{enumerate}
  \item \textbf{Stable under benign traffic.} False-alarm rate stays below 0.5\%
  without drift (0.00\% at 3~min, 0.38\% at 1~h, 0.42\% over 22{,}644 flows at 6~h),
  ruling out any unit/scaling defect (which would spike FAR toward 100\%).
  \item \textbf{Zero-shot does not transfer.} Benchmark excellence collapses on live
  traffic (MCC $\approx 0$; $-0.22$ for CIC); a three-domain analysis attributes this
  to genuine distribution shift (domain-classifier accuracy $\approx 0.99$ among
  CIC/UNSW/AWS), concentrated in rate features and backward packet size.
  \item \textbf{Cheaply repairable.} 1\% AWS labels (about 107 flows) lift MCC from
  $\approx 0$ to 0.49; 10--20\% reaches $\approx 0.59$, on par with AWS-only (0.598).
  We report MCC because AWS-test is $\approx 99\%$ attack, where F1 would be
  misleadingly high.
\end{enumerate}
A stricter episode-level split shows the few-shot gain is smaller and
episode-dependent; we present the flow-level operating point as practical, not a
universal prescription. The most robust evidence for low-cost adaptation remains the
offline cross-network experiments (Section~5).

\section{Discussion and Limitations}
The findings form a coherent narrative: single-source NIDS collapse cross-network; a
directly measured distribution shift (not feature insufficiency) is the major
contributor; and minimal calibration or mixup recovers performance, confirmed on live
AWS traffic. \textbf{Limitations:} the direction asymmetry is not fully resolved;
\begin{table}[htbp]
\centering
\caption{Domain-dispersion quantification: variance (joint $z$-space) per SFM feature.
UNSW is far more dispersed, explaining the adaptation-direction asymmetry.}
\label{tab:dispersion}
\footnotesize
\begin{tabular}{lrr}
\toprule
SFM feature & CIC var. & UNSW var. \\
\midrule
Flow Duration    & $1.11$ & $0.00$ \\
Tot Fwd Pkts     & $1.11$ & $0.01$ \\
Tot Bwd Pkts     & $1.05$ & $0.51$ \\
TotLen Fwd Pkts  & $0.47$ & $5.90$ \\
TotLen Bwd Pkts  & $1.06$ & $0.42$ \\
Fwd Pkt Len Mean & $0.45$ & $5.23$ \\
Bwd Pkt Len Mean & $0.87$ & $2.17$ \\
Flow Byts/s      & $0.00$ & $9.00$ \\
Bwd Pkts/s       & $0.01$ & $9.50$ \\
\midrule
\textbf{Mean}    & $\mathbf{0.68}$ & $\mathbf{3.64}$ \\
\bottomrule
\end{tabular}
\end{table}

using a single XGBoost is a deliberate edge-efficiency choice (2.9~MB,
sub-millisecond inference) rather than a limitation, but a head-to-head against deep
domain-invariant baselines (DI-NIDS, DANN) needs a different stack and is future work.
Adversarial robustness is deliberately scoped out: preliminary results show
adversarial training improves in-domain evasion robustness but does not close the
cross-network gap---confirming the two problems need different solutions.

\section{Conclusion and Future Work}
We presented a statistically validated cross-tool Semantic Feature Mapping and---
uniquely among cross-dataset NIDS studies---validated the full recipe on live AWS
traffic: stable false-alarm behavior, a quantified three-domain distribution shift,
and few-shot recovery at minimal labeling cost. Future work: online/streaming
continuous adaptation with drift detection; automated few-shot labeling (active
learning, SIEM/SOC integration) with a trusted-holdout guardrail; robust ensembles
under edge budgets; and multi-class attack-category extension. Adversarial robustness
is pursued in a separate study.

\section*{Acknowledgements}
The authors thank Politeknik Elektronika Negeri Surabaya (PENS) for computing
support. (Add funding/grant sources before submission if applicable.)

\section*{Nomenclature}
\begin{tabular}{ll}
NIDS & Network Intrusion Detection System \\
SFM & Semantic Feature Mapping (statistically validated cross-tool mapping) \\
MCC & Matthews Correlation Coefficient \\
FAR & False-Alarm Rate \\
IAT & Inter-Arrival Time (flow feature) \\
CORAL & CORrelation ALignment (domain adaptation) \\
DANN & Domain-Adversarial Neural Network \\
DI-NIDS & Domain-Invariant NIDS \\
$W_1$ & 1-Wasserstein distance between domain distributions \\
$\sigma,\mu$ & standard deviation and mean (z-score scaling) \\
\end{tabular}

\begin{thebibliography}{99}
\bibitem{cantone} M.~Cantone, C.~Marrocco, and A.~Bria, Machine learning in network intrusion detection: A cross-dataset generalization study, \emph{IEEE Access}, 12, pp.~144491--144509, 2024.
\bibitem{dhooge} L.~D'Hooge, T.~Wauters, B.~Volckaert, and F.~De~Turck, Inter-dataset generalization strength of supervised machine learning methods for intrusion detection, \emph{J. Inf. Secur. Appl.}, 54, art.~102564, 2020.
\bibitem{verkerken} M.~Verkerken, L.~D'Hooge, T.~Wauters, B.~Volckaert, and F.~De~Turck, Towards model generalization for intrusion detection: Unsupervised machine learning techniques, \emph{J. Netw. Syst. Manage.}, 30(1), pp.~1--25, 2022.
\bibitem{sarhan} M.~Sarhan, S.~Layeghy, and M.~Portmann, Towards a standard feature set for network intrusion detection system datasets, \emph{Mobile Netw. Appl.}, 27(1), pp.~357--370, 2022.
\bibitem{sarhannf} M.~Sarhan, S.~Layeghy, N.~Moustafa, and M.~Portmann, NetFlow datasets for machine learning-based network intrusion detection systems, in \emph{Big Data Technologies and Applications}, Springer, 2021.
\bibitem{engelen} G.~Engelen, V.~Rimmer, and W.~Joosen, Troubleshooting an intrusion detection dataset: The CICIDS2017 case study, \emph{Proc. IEEE SPW}, pp.~7--12, 2021.
\bibitem{rosay} A.~Rosay, E.~Cheval, F.~Carlier, and P.~Leroux, Network intrusion detection: A comprehensive analysis of CIC-IDS2017, \emph{Proc. 8th ICISSP}, pp.~25--36, 2022.
\bibitem{lanvin} M.~Lanvin, P.-F.~Gimenez, Y.~Han, F.~Majorczyk, L.~M\'e, and \'E.~Totel, Errors in the CICIDS2017 dataset and the significant differences in detection performances it makes, \emph{Proc. CRiSIS}, pp.~18--33, 2022.
\bibitem{liu} L.~Liu, G.~Engelen, T.~Lynar, D.~Essam, and W.~Joosen, Error prevalence in NIDS datasets: A case study on CIC-IDS-2017 and CSE-CIC-IDS-2018, \emph{Proc. IEEE CNS}, pp.~254--262, 2022.
\bibitem{catillo} M.~Catillo, A.~Pecchia, and U.~Villano, Machine learning on public intrusion datasets: Academic hype or concrete advances in NIDS?, \emph{Proc. IEEE/IFIP DSN-S}, pp.~132--136, 2023.
\bibitem{layeghy} S.~Layeghy, M.~Gallagher, and M.~Portmann, Benchmarking the benchmark---Comparing synthetic and real-world network IDS datasets, \emph{J. Inf. Secur. Appl.}, 80, art.~103689, 2024.
\bibitem{coral} B.~Sun and K.~Saenko, Deep CORAL: Correlation alignment for deep domain adaptation, \emph{Proc. ECCV Workshops}, pp.~443--450, 2016.
\bibitem{xgboost} T.~Chen and C.~Guestrin, XGBoost: A scalable tree boosting system, \emph{Proc. 22nd ACM SIGKDD}, pp.~785--794, 2016.
\bibitem{cicids} I.~Sharafaldin, A.~H.~Lashkari, and A.~A.~Ghorbani, Toward generating a new intrusion detection dataset and intrusion traffic characterization, \emph{Proc. 4th ICISSP}, pp.~108--116, 2018.
\bibitem{unsw} N.~Moustafa and J.~Slay, UNSW-NB15: A comprehensive data set for network intrusion detection systems, \emph{Proc. MilCIS}, pp.~1--6, 2015.
\bibitem{unsweval} N.~Moustafa and J.~Slay, The evaluation of network anomaly detection systems: Statistical analysis of the UNSW-NB15 data set and the comparison with the KDD99 data set, \emph{Inf. Secur. J.}, 25(1--3), pp.~18--31, 2016.
\bibitem{khraisat} A.~Khraisat, I.~Gondal, P.~Vamplew, and J.~Kamruzzaman, Survey of intrusion detection systems: Techniques, datasets and challenges, \emph{Cybersecurity}, 2(1), pp.~1--22, 2019.
\bibitem{buczak} A.~L.~Buczak and E.~Guven, A survey of data mining and machine learning methods for cyber security intrusion detection, \emph{IEEE Commun. Surveys Tuts.}, 18(2), pp.~1153--1176, 2016.
\bibitem{cremer} F.~Cremer, B.~Sheehan, M.~Fortmann, A.~N.~Kia, M.~Mullins, F.~Murphy, and S.~Materne, Cyber risk and cybersecurity: A systematic review of data availability, \emph{Geneva Papers Risk Insur.}, 47(3), pp.~698--736, 2022.
\bibitem{shone} N.~Shone, T.~N.~Ngoc, V.~D.~Phai, and Q.~Shi, A deep learning approach to network intrusion detection, \emph{IEEE Trans. Emerg. Topics Comput. Intell.}, 2(1), pp.~41--50, 2018.
\bibitem{wu} Z.~Wu, H.~Zhang, P.~Wang, and Z.~Sun, RTIDS: A robust transformer-based approach for intrusion detection system, \emph{IEEE Access}, 10, pp.~64375--64387, 2022.
\bibitem{husain} A.~Husain, A.~Salem, C.~Jim, and G.~Dimitoglou, Development of an efficient network intrusion detection model using XGBoost on the UNSW-NB15 dataset, \emph{Proc. IEEE ISSPIT}, pp.~1--7, 2019.
\end{thebibliography}

\end{document}
